\IfFileExists{stacks-project.cls}{%
\documentclass{stacks-project}
}{%
\documentclass{amsart}
}

% For dealing with references we use the comment environment
\usepackage{verbatim}
\newenvironment{reference}{\comment}{\endcomment}
%\newenvironment{reference}{}{}
\newenvironment{slogan}{\comment}{\endcomment}
\newenvironment{history}{\comment}{\endcomment}

% For commutative diagrams we use Xy-pic
\usepackage[all]{xy}

% We use 2cell for 2-commutative diagrams.
\xyoption{2cell}
\UseAllTwocells

% We use multicol for the list of chapters between chapters
\usepackage{multicol}

% This is generally recommended for better output
\usepackage{lmodern}
\usepackage[T1]{fontenc}

% For cross-file-references
\usepackage{xr-hyper}

% Package for hypertext links:
\usepackage{hyperref}

% For any local file, say "hello.tex" you want to link to please
% use \externaldocument[hello-]{hello}
\externaldocument[introduction-]{introduction}
\externaldocument[conventions-]{conventions}
\externaldocument[sets-]{sets}
\externaldocument[categories-]{categories}
\externaldocument[topology-]{topology}
\externaldocument[sheaves-]{sheaves}
\externaldocument[sites-]{sites}
\externaldocument[stacks-]{stacks}
\externaldocument[fields-]{fields}
\externaldocument[algebra-]{algebra}
\externaldocument[brauer-]{brauer}
\externaldocument[homology-]{homology}
\externaldocument[derived-]{derived}
\externaldocument[simplicial-]{simplicial}
\externaldocument[more-algebra-]{more-algebra}
\externaldocument[smoothing-]{smoothing}
\externaldocument[modules-]{modules}
\externaldocument[sites-modules-]{sites-modules}
\externaldocument[injectives-]{injectives}
\externaldocument[cohomology-]{cohomology}
\externaldocument[sites-cohomology-]{sites-cohomology}
\externaldocument[dga-]{dga}
\externaldocument[dpa-]{dpa}
\externaldocument[sdga-]{sdga}
\externaldocument[hypercovering-]{hypercovering}
\externaldocument[schemes-]{schemes}
\externaldocument[constructions-]{constructions}
\externaldocument[properties-]{properties}
\externaldocument[morphisms-]{morphisms}
\externaldocument[coherent-]{coherent}
\externaldocument[divisors-]{divisors}
\externaldocument[limits-]{limits}
\externaldocument[varieties-]{varieties}
\externaldocument[topologies-]{topologies}
\externaldocument[descent-]{descent}
\externaldocument[perfect-]{perfect}
\externaldocument[more-morphisms-]{more-morphisms}
\externaldocument[flat-]{flat}
\externaldocument[groupoids-]{groupoids}
\externaldocument[more-groupoids-]{more-groupoids}
\externaldocument[etale-]{etale}
\externaldocument[chow-]{chow}
\externaldocument[intersection-]{intersection}
\externaldocument[pic-]{pic}
\externaldocument[weil-]{weil}
\externaldocument[adequate-]{adequate}
\externaldocument[dualizing-]{dualizing}
\externaldocument[duality-]{duality}
\externaldocument[discriminant-]{discriminant}
\externaldocument[derham-]{derham}
\externaldocument[local-cohomology-]{local-cohomology}
\externaldocument[algebraization-]{algebraization}
\externaldocument[curves-]{curves}
\externaldocument[resolve-]{resolve}
\externaldocument[models-]{models}
\externaldocument[functors-]{functors}
\externaldocument[equiv-]{equiv}
\externaldocument[pione-]{pione}
\externaldocument[etale-cohomology-]{etale-cohomology}
\externaldocument[proetale-]{proetale}
\externaldocument[relative-cycles-]{relative-cycles}
\externaldocument[more-etale-]{more-etale}
\externaldocument[trace-]{trace}
\externaldocument[crystalline-]{crystalline}
\externaldocument[spaces-]{spaces}
\externaldocument[spaces-properties-]{spaces-properties}
\externaldocument[spaces-morphisms-]{spaces-morphisms}
\externaldocument[decent-spaces-]{decent-spaces}
\externaldocument[spaces-cohomology-]{spaces-cohomology}
\externaldocument[spaces-limits-]{spaces-limits}
\externaldocument[spaces-divisors-]{spaces-divisors}
\externaldocument[spaces-over-fields-]{spaces-over-fields}
\externaldocument[spaces-topologies-]{spaces-topologies}
\externaldocument[spaces-descent-]{spaces-descent}
\externaldocument[spaces-perfect-]{spaces-perfect}
\externaldocument[spaces-more-morphisms-]{spaces-more-morphisms}
\externaldocument[spaces-flat-]{spaces-flat}
\externaldocument[spaces-groupoids-]{spaces-groupoids}
\externaldocument[spaces-more-groupoids-]{spaces-more-groupoids}
\externaldocument[bootstrap-]{bootstrap}
\externaldocument[spaces-pushouts-]{spaces-pushouts}
\externaldocument[spaces-chow-]{spaces-chow}
\externaldocument[groupoids-quotients-]{groupoids-quotients}
\externaldocument[spaces-more-cohomology-]{spaces-more-cohomology}
\externaldocument[spaces-simplicial-]{spaces-simplicial}
\externaldocument[spaces-duality-]{spaces-duality}
\externaldocument[formal-spaces-]{formal-spaces}
\externaldocument[restricted-]{restricted}
\externaldocument[spaces-resolve-]{spaces-resolve}
\externaldocument[formal-defos-]{formal-defos}
\externaldocument[defos-]{defos}
\externaldocument[cotangent-]{cotangent}
\externaldocument[examples-defos-]{examples-defos}
\externaldocument[algebraic-]{algebraic}
\externaldocument[examples-stacks-]{examples-stacks}
\externaldocument[stacks-sheaves-]{stacks-sheaves}
\externaldocument[criteria-]{criteria}
\externaldocument[artin-]{artin}
\externaldocument[quot-]{quot}
\externaldocument[stacks-properties-]{stacks-properties}
\externaldocument[stacks-morphisms-]{stacks-morphisms}
\externaldocument[stacks-limits-]{stacks-limits}
\externaldocument[stacks-cohomology-]{stacks-cohomology}
\externaldocument[stacks-perfect-]{stacks-perfect}
\externaldocument[stacks-introduction-]{stacks-introduction}
\externaldocument[stacks-more-morphisms-]{stacks-more-morphisms}
\externaldocument[stacks-geometry-]{stacks-geometry}
\externaldocument[moduli-]{moduli}
\externaldocument[moduli-curves-]{moduli-curves}
\externaldocument[examples-]{examples}
\externaldocument[exercises-]{exercises}
\externaldocument[guide-]{guide}
\externaldocument[desirables-]{desirables}
\externaldocument[coding-]{coding}
\externaldocument[obsolete-]{obsolete}
\externaldocument[fdl-]{fdl}
\externaldocument[index-]{index}

% Theorem environments.
%
\theoremstyle{plain}
\newtheorem{theorem}[subsection]{Theorem}
\newtheorem{proposition}[subsection]{Proposition}
\newtheorem{lemma}[subsection]{Lemma}

\theoremstyle{definition}
\newtheorem{definition}[subsection]{Definition}
\newtheorem{example}[subsection]{Example}
\newtheorem{exercise}[subsection]{Exercise}
\newtheorem{situation}[subsection]{Situation}

\theoremstyle{remark}
\newtheorem{remark}[subsection]{Remark}
\newtheorem{remarks}[subsection]{Remarks}

\numberwithin{equation}{subsection}

% Macros
%
\def\lim{\mathop{\mathrm{lim}}\nolimits}
\def\colim{\mathop{\mathrm{colim}}\nolimits}
\def\Spec{\mathop{\mathrm{Spec}}}
\def\Hom{\mathop{\mathrm{Hom}}\nolimits}
\def\Ext{\mathop{\mathrm{Ext}}\nolimits}
\def\SheafHom{\mathop{\mathcal{H}\!\mathit{om}}\nolimits}
\def\SheafExt{\mathop{\mathcal{E}\!\mathit{xt}}\nolimits}
\def\Sch{\mathit{Sch}}
\def\Mor{\mathop{\mathrm{Mor}}\nolimits}
\def\Ob{\mathop{\mathrm{Ob}}\nolimits}
\def\Sh{\mathop{\mathit{Sh}}\nolimits}
\def\NL{\mathop{N\!L}\nolimits}
\def\CH{\mathop{\mathrm{CH}}\nolimits}
\def\proetale{{pro\text{-}\acute{e}tale}}
\def\etale{{\acute{e}tale}}
\def\QCoh{\mathit{QCoh}}
\def\Ker{\mathop{\mathrm{Ker}}}
\def\Im{\mathop{\mathrm{Im}}}
\def\Coker{\mathop{\mathrm{Coker}}}
\def\Coim{\mathop{\mathrm{Coim}}}

% Boxtimes
%
\DeclareMathSymbol{\boxtimes}{\mathbin}{AMSa}{"02}

%
% Macros for moduli stacks/spaces
%
\def\QCohstack{\mathcal{QC}\!\mathit{oh}}
\def\Cohstack{\mathcal{C}\!\mathit{oh}}
\def\Spacesstack{\mathcal{S}\!\mathit{paces}}
\def\Quotfunctor{\mathrm{Quot}}
\def\Hilbfunctor{\mathrm{Hilb}}
\def\Curvesstack{\mathcal{C}\!\mathit{urves}}
\def\Polarizedstack{\mathcal{P}\!\mathit{olarized}}
\def\Complexesstack{\mathcal{C}\!\mathit{omplexes}}
% \Pic is the operator that assigns to X its picard group, usage \Pic(X)
% \Picardstack_{X/B} denotes the Picard stack of X over B
% \Picardfunctor_{X/B} denotes the Picard functor of X over B
\def\Pic{\mathop{\mathrm{Pic}}\nolimits}
\def\Picardstack{\mathcal{P}\!\mathit{ic}}
\def\Picardfunctor{\mathrm{Pic}}
\def\Deformationcategory{\mathcal{D}\!\mathit{ef}}
\usepackage{cleveref}
\crefname{originallemma}{original lemma}{original lemmas}
\Crefname{originallemma}{Original lemma}{Original lemmas}
\makeatletter
\newlabel{homology-lemma-double-complex-gives-resolution}{{0133}{source}{}{sourcepointer-1}{}}
\newlabel{homology-lemma-double-complex-gives-resolution@cref}{{[originallemma][0][]0133}{[1][0][]source}{}{}{}}
\makeatother
\crefname{originallemma}{original lemma}{original lemmas}
\Crefname{originallemma}{Original lemma}{Original lemmas}
\makeatletter
\newlabel{homology-lemma-first-quadrant-ss}{{0132}{source}{}{sourcepointer-2}{}}
\newlabel{homology-lemma-first-quadrant-ss@cref}{{[originallemma][0][]0132}{[1][0][]source}{}{}{}}
\makeatother
\crefname{originalreference}{original reference}{original references}
\Crefname{originalreference}{Original reference}{Original references}
\makeatletter
\newlabel{lemma-describe-sheaves-simplicial-site}{{09VM}{source}{}{sourcepointer-3}{}}
\newlabel{lemma-describe-sheaves-simplicial-site@cref}{{[originalreference][0][]09VM}{[1][0][]source}{}{}{}}
\makeatother
\crefname{originalreference}{original reference}{original references}
\Crefname{originalreference}{Original reference}{Original references}
\makeatletter
\newlabel{lemma-restriction-injective-to-component}{{09W1}{source}{}{sourcepointer-4}{}}
\newlabel{lemma-restriction-injective-to-component@cref}{{[originalreference][0][]09W1}{[1][0][]source}{}{}{}}
\makeatother
\crefname{originalreference}{original reference}{original references}
\Crefname{originalreference}{Original reference}{Original references}
\makeatletter
\newlabel{lemma-restriction-to-components}{{09W0}{source}{}{sourcepointer-5}{}}
\newlabel{lemma-restriction-to-components@cref}{{[originalreference][0][]09W0}{[1][0][]source}{}{}{}}
\makeatother
\crefname{originalreference}{original reference}{original references}
\Crefname{originalreference}{Original reference}{Original references}
\makeatletter
\newlabel{lemma-simplicial-resolution-Z}{{09W5}{source}{}{sourcepointer-6}{}}
\newlabel{lemma-simplicial-resolution-Z@cref}{{[originalreference][0][]09W5}{[1][0][]source}{}{}{}}
\makeatother
\crefname{originallemma}{original lemma}{original lemmas}
\Crefname{originallemma}{Original lemma}{Original lemmas}
\makeatletter
\newlabel{sites-cohomology-lemma-injective-abelian-sheaf-injective-presheaf}{{03F6}{source}{}{sourcepointer-7}{}}
\newlabel{sites-cohomology-lemma-injective-abelian-sheaf-injective-presheaf@cref}{{[originallemma][0][]03F6}{[1][0][]source}{}{}{}}
\makeatother
\title{The E\_1 spectral sequence for a simplicial topological space with E\_1\textasciicircum{}\{p,q\} = H\textasciicircum{}q(X\_p, F\_p)}
\author{The Stacks Project Authors}
\date{stacks-project:spaces-simplicial.tex:local-snapshot-20261012; GNU FDL 1.2 or later}
\begin{document}
\maketitle
\noindent\textit{Editorial note.} All delivered mathematical text is the author's own, copied from the pinned Stacks Project chapter file sources/stacks-spaces-simplicial/pinned-original-source.tex; no statement or proof step was rewritten, reordered or removed.\par\medskip
\noindent\textit{Editorial note.} The counted proof is the author's complete argument: functoriality in F (statement line 632) rides on the cited Homology first-quadrant spectral sequence lemma, as recorded for the module-twin card L46178; this lemma is the (2d) ingredient of card L46177 (augmentation-pushforward-higher-direct-image); the identification A\textasciicircum{}\{p,q\} = Mor\_PAb(Z\_\{X\_p\}, I\textasciicircum{}q) is at source lines 637-641.\par\medskip
\noindent\textit{Editorial note.} The label lemma-simplicial-sheaf-cohomology names the counted statement, and the proof environment immediately following it in the source is delivered verbatim as its proof.\par\medskip
\noindent\textit{Editorial note.} The counted statement is carried by the author's own lemma environment in the chapter source.\par\medskip
\noindent\textit{Allowed standard prerequisites (editorial).}\begin{itemize}
\item Setup: a simplicial topological space and an abelian sheaf. Statement lines 624-630. X a simplicial topological space, F an abelian sheaf on X\_\{Zar\}; the spectral sequence is claimed with E\_1\textasciicircum{}\{p,q\} = H\textasciicircum{}q(X\_p, F\_p), functorial in F.
\item The Z-resolution on the simplicial site. Same chapter, \textbackslash{}label\{lemma-simplicial-resolution-Z\} and Cohomology on Sites, Lemma \textbackslash{}label\{sites-cohomology-lemma-injective-abelian-sheaf-injective-presheaf\}. Proof lines 634-641: A\textasciicircum{}\{p,q\} = Gamma(X\_p, I\_p\textasciicircum{}q) = Mor\_PAb(Z\_\{X\_p\}, I\textasciicircum{}q); the rows of the double complex are exact in positive degrees and evaluate to Gamma(X\_\{Zar\}, I\textasciicircum{}q) in degree 0.
\item Exactness of restriction and injectivity of the restricted sheaves. Same chapter, \textbackslash{}label\{lemma-restriction-to-components\} (card L46149) and \textbackslash{}label\{lemma-restriction-injective-to-component\}. Proof lines 641-648: restriction is exact, so F\_p -> I\_p\textasciicircum{}bullet is a resolution; the I\_p\textasciicircum{}q are injective abelian sheaves on X\_p.
\item The homological-algebra machinery. Homology, Lemmas \textbackslash{}label\{homology-lemma-first-quadrant-ss\} and \textbackslash{}label\{homology-lemma-double-complex-gives-resolution\}. Proof lines 648-651: the first-quadrant spectral sequence lemma and the double-complex resolution lemma conclude.
\end{itemize}\medskip
\ifdefined\hypertarget\hypertarget{sourcepointer-1}{}\fi
\noindent\textit{Source reference.} Original lemma 0133: sources/stacks-chow/tags.txt, source lines 1421--1421. Prerequisite pointer invoked at proof line(s) 665; it names a standard fact of a sibling chapter that the counted argument uses, and every step of the counted proof built on it is retained verbatim in the delivered body.\par\medskip
\ifdefined\hypertarget\hypertarget{sourcepointer-2}{}\fi
\noindent\textit{Source reference.} Original lemma 0132: sources/stacks-chow/tags.txt, source lines 1420--1420. Prerequisite pointer invoked at proof line(s) 664; it names a standard fact of a sibling chapter that the counted argument uses, and every step of the counted proof built on it is retained verbatim in the delivered body.\par\medskip
\ifdefined\hypertarget\hypertarget{sourcepointer-3}{}\fi
\noindent\textit{Source reference.} Original reference 09VM: sources/stacks-spaces-simplicial/pinned-original-source.tex, source lines 352--352. Prerequisite pointer invoked at proof line(s) 642; it names a standard fact of the same chapter that the counted argument uses, and every step of the counted proof built on it is retained verbatim in the delivered body.\par\medskip
\ifdefined\hypertarget\hypertarget{sourcepointer-4}{}\fi
\noindent\textit{Source reference.} Original reference 09W1: sources/stacks-spaces-simplicial/pinned-original-source.tex, source lines 492--492. Prerequisite pointer invoked at proof line(s) 661; it names a standard fact of the same chapter that the counted argument uses, and every step of the counted proof built on it is retained verbatim in the delivered body.\par\medskip
\ifdefined\hypertarget\hypertarget{sourcepointer-5}{}\fi
\noindent\textit{Source reference.} Original reference 09W0: sources/stacks-spaces-simplicial/pinned-original-source.tex, source lines 456--456. Prerequisite pointer invoked at proof line(s) 654; it names a standard fact of the same chapter that the counted argument uses, and every step of the counted proof built on it is retained verbatim in the delivered body.\par\medskip
\ifdefined\hypertarget\hypertarget{sourcepointer-6}{}\fi
\noindent\textit{Source reference.} Original reference 09W5: sources/stacks-spaces-simplicial/pinned-original-source.tex, source lines 590--590. Prerequisite pointer invoked at proof line(s) 648; it names a standard fact of the same chapter that the counted argument uses, and every step of the counted proof built on it is retained verbatim in the delivered body.\par\medskip
\ifdefined\hypertarget\hypertarget{sourcepointer-7}{}\fi
\noindent\textit{Source reference.} Original lemma 03F6: sources/stacks-chow/tags.txt, source lines 4481--4481. Prerequisite pointer invoked at proof line(s) 650; it names a standard fact of a sibling chapter that the counted argument uses, and every step of the counted proof built on it is retained verbatim in the delivered body.\par\medskip
\begin{lemma}
\label{lemma-simplicial-sheaf-cohomology}
Let $X$ be a simplicial topological space. Let $\mathcal{F}$ be an abelian
sheaf on $X$. There is a spectral sequence $(E_r, d_r)_{r \geq 0}$ with
$$
E_1^{p, q} = H^q(X_p, \mathcal{F}_p)
$$
converging to $H^{p + q}(X_{Zar}, \mathcal{F})$.
This spectral sequence is functorial in $\mathcal{F}$.
\end{lemma}

\begin{proof}
Let $\mathcal{F} \to \mathcal{I}^\bullet$ be an injective resolution.
Consider the double complex with terms
$$
A^{p, q} = \mathcal{I}^q(X_p)
$$
and first differential given by the alternating sum along the maps
$d^{p + 1}_i$-maps $\mathcal{I}_p^q \to \mathcal{I}_{p + 1}^q$, see
Lemma \ref{lemma-describe-sheaves-simplicial-site}. Note that
$$
A^{p, q} = \Gamma(X_p, \mathcal{I}_p^q) =
\Mor_{\textit{PSh}}(h_{X_p}, \mathcal{I}^q) =
\Mor_{\textit{PAb}}(\mathbf{Z}_{X_p}, \mathcal{I}^q)
$$
Hence it follows from Lemma \ref{lemma-simplicial-resolution-Z} and
Cohomology on Sites, Lemma
\ref{sites-cohomology-lemma-injective-abelian-sheaf-injective-presheaf}
that the rows of the double complex are exact in positive degrees and
evaluate to $\Gamma(X_{Zar}, \mathcal{I}^q)$ in degree $0$.
On the other hand, since restriction is exact
(Lemma \ref{lemma-restriction-to-components})
the map
$$
\mathcal{F}_p \to \mathcal{I}_p^\bullet
$$
is a resolution. The sheaves $\mathcal{I}_p^q$ are injective
abelian sheaves on $X_p$
(Lemma \ref{lemma-restriction-injective-to-component}).
Hence the cohomology of the columns computes the groups
$H^q(X_p, \mathcal{F}_p)$. We conclude by applying
Homology, Lemmas \ref{homology-lemma-first-quadrant-ss} and
\ref{homology-lemma-double-complex-gives-resolution}.
\end{proof}

\end{document}
